\subsection{Published generalized slack projection}\label{sec:gsp}
Generalized slack projection removes certain explicit slack variables from an
OSQP-style split by replacing the box projection with a proximal operation. The
published derivation and native implementation are due to Schwan, Lew and
Plancher, building on slack elimination in ADMM~\cite{gsp,lew}. The implementation
here preserves the authors' source identity and supported penalty families. It
does not equate a citation to an executed method: compatible synthetic cases
actually call the native backend and recover original-space objectives and duals.

Let $C$ be the row interval box, $\Pi_C$ its projection, and $\Phi$ the supported
slack penalty. The key residual/projection construction is
\begin{equation}\label{eq:gsp-residual}
 R_C(v)=v-\Pi_C(v),\qquad
 T(v)=\Pi_C(v)+\prox_{\Phi/\rho}\bigl(R_C(v)\bigr),\quad \rho>0,
\end{equation}
with the corrective-slack versus violation-residual convention fixed consistently.
The proximal operation itself is
\begin{equation}\label{eq:gsp-prox}
 \prox_{\Phi/\rho}(r)=\argmin_w\left\{\Phi(w)+\tfrac\rho2\|w-r\|_2^2\right\}.
\end{equation}
For the published unconditional penalties, sign symmetry and contraction behavior
permit the generalized projection to represent the eliminated slack problem.
These are mathematical assumptions, not implementation options that can be
ignored when a requested lower price differs from an upper price.

The supported published paths include symmetric quadratic, linear, elastic and
selected Huber or grouped penalties. In C13 the mixed objective exercises that
broader native penalty representation. By contrast, the authored side-asymmetric
energy variants cannot be sent to this backend as if their coefficients were
symmetric. The adapter returns an incompatible state with a reason rather than
averaging the coefficients or solving a symmetric surrogate. This refusal is
visible in both the method inventory and the cross-case displays.

The modified proximal operation changes the row update without necessarily
enlarging the original $x,z$ principal system. That property is useful when a
problem repeatedly updates its vectors while retaining matrices and penalty
structure. It does not remove the sparse system solve, guarantee invariant fill,
or imply that every numerical parameter change is free. An adaptive ADMM penalty
can require refactorization even when the matrix dimensions remain identical.
The computational comparison consequently measures the complete specified
operation rather than counting auxiliaries as a substitute for elapsed time.

Algorithm~\ref{alg:gsp} separates compatibility, numerical execution and original
acceptance. A prepared workspace may retain fixed matrices, a fixed hard mask,
the finite-side pattern and the admitted penalty family. Only permitted vector
changes are applied during repeated execution. Passing a different group layout
or a different coefficient structure is not a valid vector-only update. The
same restrictions apply to references and the candidate, avoiding a comparison
that grants reuse only to one selected implementation.

\begin{algorithm}[htbp]
\caption{Execute the published generalized projection only on supported penalty semantics, retaining the original row system and independent objective and feasibility recovery.}\label{alg:gsp}
\begin{algorithmic}[1]
\Require Original problem, published penalty family and pinned backend
\State Check symmetry, grouping and finite-side assumptions
\If{the requested objective is unsupported}
\State Return incompatible with the precise semantic reason
\EndIf
\State Prepare the original-row linear system and admitted penalty descriptors
\State Iterate the native solve and generalized row operation \eqref{eq:gsp-residual}
\State Recover original $x,y$; evaluate the requested original objective
\State Report termination, hard violation, disagreement and numerical diagnostics
\end{algorithmic}
\end{algorithm}

Convergence follows the underlying convex splitting result only within its
stated assumptions, including appropriate solution existence and numerical
parameter conditions~\cite{admm,eckstein}. The published proximal formula is not
a blanket theorem for a sign-asymmetric, variable-metric or learned-parameter
extension. The separate candidate derivations state precisely which broader
penalties and fixed metrics are considered. Their accuracy and convergence
claims do not retroactively enlarge the domain of the authors' code.

Source attribution is substantive. The adapted native code is distributed by
the separate engine; nested OSQP and QDLDL terms are retained alongside the
OSQPSlack notice. The building source is reproduced with its original mixed
penalties before authored variants are introduced. Preserving those distinctions
allows a reader to identify what is reproduced, what is an adapter, and what is
a mathematical extension. A plot that omits incompatible cells must retain their
reasons elsewhere; otherwise absence can be mistaken for a successful zero-cost
solution or evidence that a supported reference was never executed.

The common projection notation does not erase differences between separable
and shared-group costs. With independent quadratic prices, each row operation
can be evaluated separately. A shared infinity-norm price requires a joint
operation because changing one row's largest departure can alter the group's
active coordinates. The published native descriptors encode that grouping.
The source-faithful building case exercises precisely this distinction. A
consumer that merely checks the number of rows, while dropping the group
descriptor, could accept a well-formed numerical vector for the wrong
optimization problem. Original objective reconstruction is therefore part
of the reference comparison, not an optional display calculation.
